\textbf{Pohyblivost} udává střední driftovou rychlost nosiče připadající na jednotku intenzity elektrického pole:
\[\boxed{\mu=\frac{|v_d|}{|E|}}\quad[\mathrm{m^2/(V\,s)}].\]
V kovu ji ovlivňuje rozptyl elektronů na tepelných kmitech mřížky, nečistotách, poruchách a hranicích zrn; větší rozptyl pohyblivost snižuje. Závisí tedy na teplotě, čistotě a struktuře materiálu. V kapalinách pohyblivost iontů klesá s viskozitou, přibližně $\mu\sim1/\eta$.
